\subsection{Zusammenfassung der Ergebnisse}
\label{subsec:zusammenfassung-ergebnisse}

Das Master-Template begegnet wiederholtem Setup-Aufwand und divergierenden
Projektbasen mit einem gemeinsamen Kern und deklarativer Variabilität. Es
bündelt Generierung, Konfiguration, Verantwortungsgrenzen, Arbeitsabläufe und
Wartung.

Fünf Profile kombinieren gemeinsames Vite-Frontend, profilabhängiges
FastAPI-Backend und optionale Tauri-Schicht; PostgreSQL ist nur mit Backend
wählbar. \texttt{control.py} vereinheitlicht den Entwicklungsweg. Generierte
Projekte bleiben eigenständig und erhalten keine automatischen Master-Updates.

Die lokale Abnahme eines früheren Commits stützt Profil-, PostgreSQL-, Test-,
Web-Build- und Linux-Paketierungspfade. Die aktuelle externe CI ist nur
teilweise erfolgreich; weitere Nachweise fehlen oder sind produktspezifisch.
Belegt ist damit eine kontrolliert variable Baseline, nicht Produktionsreife
oder empirische Kundeneignung.
